\chapter{Simulation and Analysis}\label{ch:simulation-and-analysis}

\section{Objectives and Status}\label{objectives-and-status}

Simulation serves four purposes: verify the collinear kinematic signs
and rank; test the nonlinear balance equations; select an initial
controller without risking hardware; and identify which parameters most
strongly affect performance. A MATLAB script named
simulate\_collinear\_mecanum\_robot.m was developed in prior project
work as a draft artifact. It used four mechanical coordinates
q=[X,Y,\ensuremath{\psi},\ensuremath{\theta}], eight mechanical states, four actuator-lag states, a
nonlinear Euler-Lagrange model, an LQR controller, sample-and-hold
control, fourth-order Runge-Kutta integration, saturation, slew limits,
and DISARMED\slash{}BALANCING\slash{}FALLEN logic. That script was created in a prior
sandbox and was not confirmed in the repository snapshot; any numerical
examples from it are not thesis results.

\begin{longtable}[]{@{}
  >{\raggedright\arraybackslash}p{(\columnwidth - 0\tabcolsep) * \real{1.0000}}@{}}
\toprule\noalign{}
\begin{minipage}[b]{\linewidth}\raggedright
\textless INSERT FIGURE 9.1: Show parameter loading, geometry\slash{}Jacobian
construction, rank checks, nonlinear plant, actuator lag and saturation,
sampled controller, state machine, RK4 integration, logging, metrics,
plots, and animation\textgreater{}\\
What the figure should show: Show parameter loading, geometry\slash{}Jacobian
construction, rank checks, nonlinear plant, actuator lag and saturation,
sampled controller, state machine, RK4 integration, logging, metrics,
plots, and animation\\
Likely source: final MATLAB simulation code\strut
\end{minipage} \\
\midrule\noalign{}
\endhead
\bottomrule\noalign{}
\endlastfoot
\end{longtable}

\begin{center}\singlespacing
\phantomsection\addcontentsline{lof}{figure}{\protect\numberline{9.1}Simulation model and numerical-integration workflow.}
\textbf{Figure 9.1. Simulation model and numerical-integration workflow.}
\end{center}

\section{Model Structure and Numerical Integration}\label{model-structure-and-numerical-integration}

The simulation state is x\_s=[X,Y,\ensuremath{\psi},\ensuremath{\theta},\ensuremath{\dot{X}},\ensuremath{\dot{Y}},\ensuremath{\dot{\psi}},\ensuremath{\dot{\theta}},\ensuremath{\tau}\_1,...,\ensuremath{\tau}\_4]\ensuremath{^{T}} when
first-order actuator torque states are retained. Parameters are defined
at the beginning of the script with units and source fields. The
controller executes at its discrete sample period while the nonlinear
plant is integrated at a smaller step. A zero-order hold preserves the
last command between controller updates.

Fourth-order Runge-Kutta integration is represented by Equation (9.1).
Step-size convergence is checked by repeating each baseline case at half
the integration step and comparing key metrics. Solver success alone is
not evidence of model correctness; dimensional checks, energy checks for
the unforced plant, zero-input equilibrium, forward\slash{}inverse kinematic
round trips, and sign tests are performed before interpreting a
response.

\begin{equation}
\mathbf{x}_{k+1}=\mathbf{x}_{k}+\frac{\Delta t}{6}
\left(\mathbf{k}_1+2\mathbf{k}_2+2\mathbf{k}_3+\mathbf{k}_4\right)
\label{eq:9-1}
\end{equation}

\begin{longtable}[]{@{}
  >{\raggedright\arraybackslash}p{(\columnwidth - 0\tabcolsep) * \real{1.0000}}@{}}
\toprule\noalign{}
\begin{minipage}[b]{\linewidth}\raggedright
\textless INSERT FIGURE 9.2: Render the robot body, four collinear wheel
centers, individual wheel rotation, roller handedness, inertial and body
axes, center of mass, pitch lean, yaw, and ground trajectory. Display
time and operating state without implying unmeasured physical
geometry\textgreater{}\\
What the figure should show: Render the robot body, four collinear wheel
centers, individual wheel rotation, roller handedness, inertial and body
axes, center of mass, pitch lean, yaw, and ground trajectory. Display
time and operating state without implying unmeasured physical geometry\\
Likely source: MATLAB visualizer using measured CAD geometry\strut
\end{minipage} \\
\midrule\noalign{}
\endhead
\bottomrule\noalign{}
\endlastfoot
\end{longtable}

\begin{center}\singlespacing
\phantomsection\addcontentsline{lof}{figure}{\protect\numberline{9.2}Three-dimensional visualization of the collinear Mecanum robot model.}
\textbf{Figure 9.2. Three-dimensional visualization of the collinear Mecanum robot model.}
\end{center}

\section{Kinematic Validation}\label{kinematic-validation}

Kinematic tests begin with unit commands in v\_x, v\_y, and \ensuremath{\dot{\psi}}. Equation
(6.3) generates wheel rates, and Equation (6.4) reconstructs the body
twist. The residual norm, H condition number, and sign pattern are
reported. Combined commands then test superposition. These tests
validate the algebra only; they do not demonstrate traction or
wheel-ground behavior.

\par\medskip\noindent\singlespacing
\phantomsection\addcontentsline{lot}{table}{\protect\numberline{9.1}Simulation cases and required outputs}
\textbf{Table 9.1. Simulation cases and required outputs}\par

\begin{longtable}[]{@{}
  >{\raggedright\arraybackslash}p{(\columnwidth - 6\tabcolsep) * \real{0.1527}}
  >{\raggedright\arraybackslash}p{(\columnwidth - 6\tabcolsep) * \real{0.2595}}
  >{\raggedright\arraybackslash}p{(\columnwidth - 6\tabcolsep) * \real{0.3130}}
  >{\raggedright\arraybackslash}p{(\columnwidth - 6\tabcolsep) * \real{0.2748}}@{}}
\toprule\noalign{}
\begin{minipage}[b]{\linewidth}\raggedright
\textbf{Case}
\end{minipage} & \begin{minipage}[b]{\linewidth}\raggedright
\textbf{Initial condition\slash{}input}
\end{minipage} & \begin{minipage}[b]{\linewidth}\raggedright
\textbf{Signals and plot}
\end{minipage} & \begin{minipage}[b]{\linewidth}\raggedright
\textbf{Derived metric}
\end{minipage} \\
\midrule\noalign{}
\endhead
\bottomrule\noalign{}
\endlastfoot
K1 longitudinal & Unit +v\_x; v\_y=\ensuremath{\dot{\psi}}=0 & Four wheel rates; reconstructed
twist & Round-trip error and signs \\
K2 lateral & Unit +v\_y; v\_x=\ensuremath{\dot{\psi}}=0 & Four wheel rates; reconstructed
twist & Round-trip error and roller pattern \\
K3 yaw & Unit +\ensuremath{\dot{\psi}}; v\_x=v\_y=0 & Wheel rates versus y\_i & Round-trip
error and yaw leverage \\
K4 combined & Bounded multiaxis command & Requested\slash{}reconstructed twist
& Maximum and RMS error \\
D1 small pitch & Small initial \ensuremath{\theta}; zero motion reference & \ensuremath{\theta}, \ensuremath{\dot{\theta}}, x,
command, state & Stability, settling, saturation \\
D2 disturbance & Defined impulse or initial angle & Pitch and command &
Recoverable disturbance\slash{}angle \\
D3 trajectory & Bounded x-y-yaw reference & Pose, velocities, pitch,
wheel commands & Tracking error and torque reserve \\
S1 sensitivity & Parameter and delay sweeps & Metric versus parameter &
Robustness margins\slash{}trends \\
\end{longtable}

\begin{longtable}[]{@{}
  >{\raggedright\arraybackslash}p{(\columnwidth - 0\tabcolsep) * \real{1.0000}}@{}}
\toprule\noalign{}
\begin{minipage}[b]{\linewidth}\raggedright
\textless INSERT FIGURE 9.3: Create three panels for pure longitudinal,
pure lateral, and pure yaw commands. Horizontal axis: time in seconds.
Vertical axes: commanded and reconstructed body velocities in m\slash{}s and
yaw rate in rad\slash{}s; include four wheel speeds in rad\slash{}s. Use the
controller sample interval and at least ten samples per fastest
displayed transient\textgreater{}\\
What the figure should show: Create three panels for pure longitudinal,
pure lateral, and pure yaw commands. Horizontal axis: time in seconds.
Vertical axes: commanded and reconstructed body velocities in m\slash{}s and
yaw rate in rad\slash{}s; include four wheel speeds in rad\slash{}s. Use the
controller sample interval and at least ten samples per fastest
displayed transient\\
Likely source: MATLAB kinematic validation logs\strut
\end{minipage} \\
\midrule\noalign{}
\endhead
\bottomrule\noalign{}
\endlastfoot
\end{longtable}

\begin{center}\singlespacing
\phantomsection\addcontentsline{lof}{figure}{\protect\numberline{9.3}Kinematic validation for independent longitudinal, lateral, and yaw commands.}
\textbf{Figure 9.3. Kinematic validation for independent longitudinal, lateral, and yaw commands.}
\end{center}

\section{Dynamic and Controller Cases}\label{dynamic-and-controller-cases}

The first dynamic case starts within the linear region with no
planar-motion request. The second applies a documented initial pitch or
torque disturbance. The third adds low-amplitude longitudinal commands,
followed by lateral and yaw commands after balance remains stable.
Saturation, state transitions, and actuator states appear on every plot
so that apparent tracking is not mistaken for unconstrained performance.

\begin{longtable}[]{@{}
  >{\raggedright\arraybackslash}p{(\columnwidth - 0\tabcolsep) * \real{1.0000}}@{}}
\toprule\noalign{}
\begin{minipage}[b]{\linewidth}\raggedright
\textless INSERT FIGURE 9.4: Plot time in seconds on the horizontal
axis. Use vertically aligned panels for pitch angle in degrees, pitch
rate in degrees\slash{}s, longitudinal position\slash{}velocity, collective command in
N\ensuremath{\cdot}m or mode-specific units, individual wheel commands, saturation flag,
and operating state. Compare nonlinear and linear predictions if both
are simulated\textgreater{}\\
What the figure should show: Plot time in seconds on the horizontal
axis. Use vertically aligned panels for pitch angle in degrees, pitch
rate in degrees\slash{}s, longitudinal position\slash{}velocity, collective command in
N\ensuremath{\cdot}m or mode-specific units, individual wheel commands, saturation flag,
and operating state. Compare nonlinear and linear predictions if both
are simulated\\
Likely source: MATLAB simulation output from the final parameter
set\strut
\end{minipage} \\
\midrule\noalign{}
\endhead
\bottomrule\noalign{}
\endlastfoot
\end{longtable}

\begin{center}\singlespacing
\phantomsection\addcontentsline{lof}{figure}{\protect\numberline{9.4}Simulated upright-balance response for the selected controller.}
\textbf{Figure 9.4. Simulated upright-balance response for the selected controller.}
\end{center}

\section{Sensitivity and Uncertainty}\label{sensitivity-and-uncertainty}

Sensitivity cases vary assembled mass, center-of-mass height, pitch
inertia, effective radius, actuator delay, torque limit, feedback age,
damping, and roller-force effectiveness over documented ranges.
One-factor sweeps show local trends; selected combinations expose
interactions. The ranges should come from measurement uncertainty or
credible configuration bounds, not arbitrary percentages.

\begin{longtable}[]{@{}
  >{\raggedright\arraybackslash}p{(\columnwidth - 0\tabcolsep) * \real{1.0000}}@{}}
\toprule\noalign{}
\begin{minipage}[b]{\linewidth}\raggedright
\textless INSERT FIGURE 9.5: Use an x-axis of the varied parameter with
units and y-axes for RMS pitch error, peak pitch, settling time, maximum
torque\slash{}current, and stability\slash{}fall outcome. Include the nominal case and
uncertainty bounds; compare delay, COM height, inertia, torque limit,
and lateral-force effectiveness\textgreater{}\\
What the figure should show: Use an x-axis of the varied parameter with
units and y-axes for RMS pitch error, peak pitch, settling time, maximum
torque\slash{}current, and stability\slash{}fall outcome. Include the nominal case and
uncertainty bounds; compare delay, COM height, inertia, torque limit,
and lateral-force effectiveness\\
Likely source: MATLAB parameter-sweep results\strut
\end{minipage} \\
\midrule\noalign{}
\endhead
\bottomrule\noalign{}
\endlastfoot
\end{longtable}

\begin{center}\singlespacing
\phantomsection\addcontentsline{lof}{figure}{\protect\numberline{9.5}Sensitivity of balance response to mass properties, delay, and actuator limits.}
\textbf{Figure 9.5. Sensitivity of balance response to mass properties, delay, and actuator limits.}
\end{center}

\clearpage
\begin{landscape}
\par\medskip\noindent\singlespacing
\phantomsection\addcontentsline{lot}{table}{\protect\numberline{9.2}Simulation parameter record}
\textbf{Table 9.2. Simulation parameter record}\par

\begin{longtable}[]{@{}
  >{\raggedright\arraybackslash}p{(\columnwidth - 8\tabcolsep) * \real{0.2482}}
  >{\raggedright\arraybackslash}p{(\columnwidth - 8\tabcolsep) * \real{0.1825}}
  >{\raggedright\arraybackslash}p{(\columnwidth - 8\tabcolsep) * \real{0.1460}}
  >{\raggedright\arraybackslash}p{(\columnwidth - 8\tabcolsep) * \real{0.2628}}
  >{\raggedright\arraybackslash}p{(\columnwidth - 8\tabcolsep) * \real{0.1606}}@{}}
\toprule\noalign{}
\begin{minipage}[b]{\linewidth}\raggedright
\textbf{Parameter}
\end{minipage} & \begin{minipage}[b]{\linewidth}\raggedright
\textbf{Value}
\end{minipage} & \begin{minipage}[b]{\linewidth}\raggedright
\textbf{Unit}
\end{minipage} & \begin{minipage}[b]{\linewidth}\raggedright
\textbf{Source\slash{}revision}
\end{minipage} & \begin{minipage}[b]{\linewidth}\raggedright
\textbf{Status}
\end{minipage} \\
\midrule\noalign{}
\endhead
\bottomrule\noalign{}
\endlastfoot
Wheel radius & \textless INSERT\textgreater{} & m & Loaded rollout\slash{}CAD &
Required \\
Wheel coordinates and signs & \textless INSERT\textgreater{} & m / --- &
Final assembly & Required \\
Body and total mass & \textless INSERT\textgreater{} & kg & Scale\slash{}CAD &
Required \\
COM height and offsets & \textless INSERT\textgreater{} & m &
CAD\slash{}physical test & Required \\
Pitch and yaw inertias & \textless INSERT\textgreater{} & kg\ensuremath{\cdot}m\ensuremath{^{2}} &
CAD\slash{}identification & Required \\
Actuator response and delay & \textless INSERT\textgreater{} & s & Bench
step test & Required \\
Torque\slash{}current\slash{}speed limits & \textless INSERT\textgreater{} & N\ensuremath{\cdot}m / A /
rad\slash{}s & Installed settings & Required \\
Controller period & \textless INSERT\textgreater{} & s & Firmware
measurement & Required \\
Integrator step & \textless INSERT\textgreater{} & s & Convergence study
& Required \\
Initial conditions & \textless INSERT\textgreater{} & SI units & Case
definition & Required \\
\end{longtable}
\end{landscape}
\clearpage

\section{Model-to-Hardware Comparison}\label{model-to-hardware-comparison}

After hardware tests, simulation and measured data should share the same
reference commands, initial conditions, sampling grid, sign convention,
and units. Compare pitch, pitch rate, wheel or actuator feedback, body
velocity, and command. Quantify time alignment and filtering. Fit only
parameters designated for identification, then validate the fitted model
on different trials. Reporting only a visually similar trace would not
establish predictive validity.

\begin{longtable}[]{@{}
  >{\raggedright\arraybackslash}p{(\columnwidth - 0\tabcolsep) * \real{1.0000}}@{}}
\toprule\noalign{}
\begin{minipage}[b]{\linewidth}\raggedright
\textless INSERT DATA: Provide final MATLAB code, parameter file, raw
simulation logs, software release, solver settings, sample interval,
integration step, and reproducibility instructions. Replace all example
values with measured or cited parameters.\textgreater{}
\end{minipage} \\
\midrule\noalign{}
\endhead
\bottomrule\noalign{}
\endlastfoot
\end{longtable}
